\chapter{Hauptteil}

Hier kommt der Hauptteil der Arbeit hinein.

\noindent
Als Platzhalter werde ich ein paar Befehle einf\"ugen, die vielleicht interessant sind...

\noindent
\textbf{Es ist selbstverst\"andlich, dass diese Zeilen aus der Datei gel\"oscht werden sollen...}

Um das Inhaltsverzeichnis zu generieren, muss $3\times$ \"ubersetzt werden! 

\section{Wie werden Umgebungen benutzt}
\begin{Satz}[Kaffeetheorem]
Ein Mathematiker ist eine Maschine, die Kaffee in Theoreme verwandelt. 
\end{Satz}

\begin{Beweis}
trivial
\end{Beweis}

\subsection{Gleichungen}
Ok, nehmen wir an, wir sollen Gleichungen ohne Nummerierung untereinander einr\"ucken. Das geht mit {\&} so

\begin{align*}
\left(a+b\right)^2&= a^2+2ab+b^2\\
&= (b+a)^2
\end{align*}


oder nummeriert

\begin{align}
\left(a+b\right)^2&=a^2+2ab+b^2\\
&= (b+a)^2
\end{align}

oder nur mit einer Nummer zentriert

\begin{align}
\begin{aligned}
\left(a+b\right)^2&= a^2+2ab+b^2\\
&=(b+a)^2
\end{aligned}
\end{align}

bzw. mit Nummer nur f\"ur die letzte Zeile

\begin{align}
\left(a+b\right)^2&= a^2+2ab+b^2\nonumber\\
&=(b+a)^2
\end{align}

Es ist auch m\"oglich Gleichungen zu labeln um sich sp\"ater im Text richtig auf sie beziehen zu k\"onnen.

\begin{align}
\label{eq:binom}
\left(a+b\right)^2= a^2+2ab+b^2
\end{align}

Die Gleichung~(\ref{eq:binom}) ist die erste binomische Formel. 


Referenzieren bzw. zitieren aus Quellen funktioniert durch \cite{ToselliWidlund}. Es sollte allerdings \textbf{immer} mit Seitenzahl oder Angabe des Theorems zitiert werden, das hei\ss t \cite[S. 127]{ToselliWidlund} oder \cite[Theorem~4.19]{ToselliWidlund}.
F\"ur die \"Ubersetzung k\"onnen folgende Kommandos hintereinander ausgef\"uhrt werden:

\begin{enumerate}
\item latex layout
\item bibtex layout
\item latex layout
\item latex layout
\item latex layout
\end{enumerate}

Es ist durchaus sinnvoll gesch\"utzte Leerzeichen zu verwenden um zwischen Theoremen und deren Nummern einen Zeilenumbruch zu vermeiden... gesch\"utztes Leerzeichen ist die Tilde...

Eine Tabelle:

\begin{table}
\small
	\begin{center}
	\begin{tabular}[c]{|c|c|c|c|c|c|c|c|c|c|c|c|}
	\hline
	\multicolumn{12}{|c|}{\textbf{A}}\\
	\hline
	&\multicolumn{5}{c|}{\textit{B}}&\multicolumn{3}{c|}{\textit{C}}&\multicolumn{3}{c|}{\textit{D}}\\
	\hline$a$& &$b$&c&$|d|$ &$\#e$                &$f$  &$g$&h&$|z|$&as&M\\
	\hline
	2& b & 32 & 153 & 264 & 0\% &1 & 213& 2&5 &6 &3\\ 
	2& c &3255 & 118 & 57 & 20\% & 123& 2&4 &4 &5 &4\\ 
	\hline
	\end{tabular}
\caption{Tabellen-Caption.}
	\label{tab:reg2}
	\end{center}
\end{table}

\begin{figure}
\includegraphics[scale=0.25]{uzk_logo}
\includegraphics[scale=0.25]{uzk_logo_grau}
\caption{Figure-Caption mit 2 Bildern.}
\end{figure}

